\begin{table}[H]
\centering
\small
\begin{tabular}{lcccc}
\toprule
model & p99.9 $\hat\delta$ & excess & $z$ & null reps above real \\
\midrule
ViT-T & $0.100$ & $+0.0069$ & $+1.3$ & 1/20 \\
ViT-S & $0.079$ & $-0.0048$ & $-1.1$ & 17/20 \\
ViT-B & $0.081$ & $-0.0065$ & $-1.8$ & 20/20 \\
ViT-L & $0.069$ & $-0.0139$ & $-4.3$ & 20/20 \\
\midrule
DINO-B & $0.054$ & $-0.0247$ & $-10.9$ & 20/20 \\
DINOv2-S & $0.053$ & $-0.0220$ & $-16.1$ & 20/20 \\
DINOv2-B & $0.035$ & $-0.0256$ & $-33.1$ & 20/20 \\
DINOv2-L & $0.032$ & $-0.0235$ & $-42.2$ & 20/20 \\
DINOv2-G & $0.032$ & $-0.0203$ & $-37.9$ & 20/20 \\
\midrule
CLIP-B & $0.087$ & $-0.0061$ & $-1.6$ & 19/20 \\
CLIP-L & $0.076$ & $-0.0114$ & $-3.1$ & 20/20 \\
SigLIP-B & $0.076$ & $-0.0090$ & $-2.7$ & 20/20 \\
\bottomrule
\end{tabular}
\caption{The ImageNet census under the supremum-robust statistic: 99.9th percentile of the four-point defect, 20 spectrum-null replicates, with the percentile rank (how many replicates exceed the real value; 20/20 = below every replicate). The ordering shifts relative to the supremum: DINO/DINOv2 carry the most budget-robust excess, while ViT-T's supremum excess does not survive ($+0.007$, above 19/20 replicates), consistent with the curvature-sign ordering of Table~\ref{tab:a10}. Large $|z|$ should be read as ``below every replicate'', not as a Gaussian tail probability.}
\label{tab:b18-p999}
\end{table}
